% c02 —— §2 辛流形（续）+ §3 分析力学，p006–p008
% 本块开头续 p5 末 定義2.1 的正文（该 definition 环境由 c01.tex 开、在此闭合），
% tag 共 13 条：(2.2)–(2.10)、(3.1)–(3.4)。

为单射（从而为同构）时，称 $\omega$ 非退化（non-degenerate）。
\end{definition}

\begin{definition}
向量空间 $V$ 上的非退化 2 形式（2 阶完全反对称张量）
\begin{equation}
    \omega\in\bigwedge^2V^*\subset V^*\otimes V^* \tag{2.2}
\end{equation}
称为 $V$ 上的辛形式（symplectic form），组 $(V,\omega)$ 称为辛向量空间（symplectic vector space）。
\end{definition}

$\omega$ 的反对称性与非退化性，可以看作是对 $J\in\End(V)\cong V^*\otimes V$ 的条件 $J^2=-\id_V$ 的类比。

\begin{exercise}
当 $\omega\in\bigwedge^2V^*$ 非退化时，证明：$V$ 是偶数维的，并且适当选取其基，可使线性映射 (2.1) 有矩阵表示
\begin{equation}
    \begin{pmatrix} 0 & -I \\ I & 0 \end{pmatrix} \tag{2.3}
\end{equation}
\end{exercise}

\begin{definition}
流形 $M$ 上的辛形式（symplectic form），是指非退化且闭的 2 次微分形式 $\omega\in\Gamma\left(\bigwedge^2T^*M\right)$。流形与其上的辛形式的组 $(M,\omega)$ 称为辛流形（symplectic manifold）。
\end{definition}

辛流形之间的「辛映射」$f\colon(M,\omega_M)\to(N,\omega_N)$ 这一概念，可以定义为满足 $f^*\omega_N=\omega_M$ 的映射 $f\colon M\to N$，但这并不是全纯映射的良好类比。例如下述结论成立：

\begin{exercise}
设 $(M,\omega_M)$ 与 $(N,\omega_N)$ 是辛流形，流形间的映射 $f\colon M\to N$ 满足 $f^*\omega_N=\omega_M$。此时证明 $\dim M\le\dim N$。
\end{exercise}

用后述的 Lagrange 子流形概念，把直积 $(M\times N,-\pi_M^*\omega_M+\pi_N^*\omega_N)$ 的 Lagrange 子流形视为辛流形之间的「映射」，在更多场合比考虑上述意义上的「辛映射」更为自然。不过，直积的 Lagrange 子流形与其说是全纯映射的类比，不如说是复几何中的对应（correspondence，即直积的复子流形）在辛几何中的类比。

另一方面，辛流形的同构可以自然地定义：

\begin{definition}
对辛流形 $(M,\omega_M)$ 与 $(N,\omega_N)$，满足 $\phi^*\omega_N=\omega_M$ 的微分同胚 $\phi\colon M\to N$ 称为辛同胚（symplectomorphism）。
\end{definition}

流形上的非退化 2 次微分形式 $\omega$ 有时也称为近辛形式。注意：辛同胚的定义中并没有用到 $\omega$ 是闭形式这一点，因此对于近辛流形之间的映射，辛同胚这一概念同样有意义。

对近辛形式而言，Nijenhuis 张量的类比物是 3 次微分形式 $d\omega$；而 Newlander–Nirenberg 定理的类比物则是 Darboux 定理。

\begin{exercise}[Darboux 定理的逆]
若近辛流形 $(M,\omega)$ 满足「对任意 $x\in M$，存在 $M$ 中包含 $x$ 的开集 $U$，使得 $(U,\omega|_U)$ 与辛向量空间的开集辛同胚」这一条件，则证明 $\omega$ 是辛形式（即 $d\omega=0$）。
\end{exercise}

\begin{exercise}[Darboux 定理]
证明：对任意辛流形 $(M,\omega)$ 与任意 $x\in M$，存在 $M$ 中包含 $x$ 的开集 $U$，使得 $(U,\omega|_U)$ 与辛向量空间的开集辛同胚。
\end{exercise}

作为包含辛结构在内的更广的概念，有 Poisson 结构。

\begin{definition}
流形 $M$ 上的 Poisson 括号（Poisson bracket），是指 $\R$ 上的双线性映射
\begin{equation}
    \{-,-\}\colon C^\infty(M)\times C^\infty(M)\to C^\infty(M) \tag{2.4}
\end{equation}
且满足反对称性
\begin{equation}
    \{f,g\}=-\{g,f\} \tag{2.5}
\end{equation}
Jacobi 恒等式
\begin{equation}
    \{f,\{g,h\}\}+\{g,\{h,f\}\}+\{h,\{f,g\}\}=0 \tag{2.6}
\end{equation}
以及 Leibniz 法则
\begin{equation}
    \{f,gh\}=\{f,g\}h+g\{f,h\} \tag{2.7}
\end{equation}
的映射。
\end{definition}

\begin{remark}
Poisson 括号的 $\R$ 双线性性、反对称性 (2.5) 以及 Jacobi 恒等式 (2.6)，可以换言之说成：Poisson 括号在 $C^\infty(M)$ 上确定了 $\R$ 上的 Lie 代数结构。
\end{remark}

\begin{remark}
Poisson 括号的 Jacobi 恒等式 (2.6)，也可以看作 Leibniz 法则
\begin{equation}
    \{f,\{g,h\}\}=\{\{f,g\},h\}+\{g,\{f,h\}\}. \tag{2.8}
\end{equation}
\end{remark}

Poisson 括号的 Leibniz 法则 (2.7)，可以换言之说成：Poisson 括号是给函数确定向量场的映射
\begin{equation}
    C^\infty(M)\to\Gamma(TM),\quad f\mapsto X_f:=\{f,-\}. \tag{2.9}
\end{equation}
把这一点与 Poisson 括号的反对称性 (2.5) 合起来看，就可以看出 Poisson 括号由 Poisson 双向量（Poisson bivector）
\begin{equation}
    \beta\in\Gamma\left(\bigwedge^2TM\right) \tag{2.10}
\end{equation}
决定。

\begin{definition}
当 Poisson 双向量非退化时，称 Poisson 括号非退化。
\end{definition}

\begin{exercise}
证明：非退化的 Poisson 括号与辛形式是等价的概念。
\end{exercise}

流形与其上的 Poisson 括号的组，称为 Poisson 流形（Poisson manifold）。

\begin{exercise}
证明：任意 Poisson 流形自然地具有这样的叶状结构，使其叶带有辛结构。
\end{exercise}

对 $H\in C^\infty(M)$，向量场 $X_H$ 称为 $H$ 确定的 Hamilton 向量场（Hamiltonian vector field）。

\begin{exercise}
证明：Poisson 流形上的 Hamilton 向量场，与习题2.6 所给的叶状结构的叶相切。
\end{exercise}

时间演化由 Hamilton 向量场给出的动力系统称为 Hamilton 系统。由习题2.7 可知，考虑 Hamilton 系统的舞台局部上可以认为是辛流形。

\section{分析力学}

Lagrange 形式的分析力学，把运动方程改写为对作用泛函（action functional）的变分原理的形式，而作用泛函由 Lagrange 函数（Lagrangian）的积分给出。

作为例子，考虑三维空间中由具有势 $V$ 的力驱动的质点的运动。设质点质量为 1，$t$ 为时刻，$\bm{x}(t)=(x(t),y(t),z(t))$ 为质点在直角坐标系中表示的运动，则 Newton 运动方程为
\begin{equation}
    \frac{d^2\bm{x}}{dt^2}=-\nabla V,\quad \nabla V=\left(\frac{\partial V}{\partial x},\frac{\partial V}{\partial y},\frac{\partial V}{\partial z}\right). \tag{3.1}
\end{equation}

此情形下的 Lagrange 函数，用动能 $K=\frac{1}{2}|\dot{\bm{x}}|^2$ 与势 $V$ 表示为 $L=K-V$。

对广义坐标 $\bm{q}=(q_1,q_2,q_3)$，质点的运动由最小作用原理确定，即使作用
\begin{equation}
    S=\int L(\bm{q},\dot{\bm{q}})dt \tag{3.2}
\end{equation}
取最小值，并满足 Euler–Lagrange 方程
% [Tutor Note: 原文 p8 此处作「Euler–Langrange 方程式」，疑为 Euler–Lagrange 之误；中文按通行拼写写作 Euler–Lagrange，请审校确认。]
\begin{equation}
    \frac{\partial L}{\partial q_i}-\frac{d}{dt}\left(\frac{\partial L}{\partial\dot{q}_i}\right)=0,\quad i=1,2,3. \tag{3.3}
\end{equation}

与 Newton 运动方程相比，Lagrange 形式的分析力学的最大优点是：作为广义坐标 $\bm{q}$，并不一定需要取直角坐标。例如当势 $V$ 为中心力、所考虑的体系球对称时，取球坐标比取直角坐标要容易看清得多。当然，可以先写出 Newton 运动方程，再用微分的链式法则把同一方程改写成球坐标形式；但与其如此，从一开始就以对称性显然可见的形式来记述，无论在理论上还是在计算上都更有利。

Lagrange 形式与 Newton 运动方程一样，运动由 2 阶常微分方程描述，方程的通解中出现变量数的 2 倍个积分常数。换一种说法：要描述某时刻体系的状态，不仅需要该时刻的广义坐标，还需要其微分（广义速度）。把位形空间（configuration space）记作 $N$，则运动方程的初始值空间就是 $N$ 的切丛 $TN$。

对作为广义坐标与广义速度的函数的 Lagrange 函数作 Legendre 变换，就得到作为广义坐标 $\bm{q}$ 与广义动量 $\bm{p}$ 的函数的 Hamilton 函数：
\begin{equation}
    p_i=\frac{\partial L}{\partial\dot{q}_i},\quad H=\sum_i\dot{q}_ip_i-L. \tag{3.4}
\end{equation}

% CHUNK c02 DONE pages p006–p008
